\documentclass[UTF8,a4paper,11pt,fontset=fandol]{ctexart}
\usepackage[margin=24mm]{geometry}
\usepackage{amsmath,amssymb,booktabs,tabularx,longtable,array}
\usepackage{xcolor,listings,xurl}
\usepackage[unicode,colorlinks=true,linkcolor=blue!45!black,urlcolor=blue!45!black]{hyperref}
\hypersetup{pdftitle={UnifiedInterpreter：HFusion/HIVM Attention 实现说明},
  pdfauthor={UnifiedInterpreter}}
\setlength{\parskip}{0.35em}
\setlength{\emergencystretch}{3em}
\renewcommand{\arraystretch}{1.25}
\newcolumntype{Y}{>{\raggedright\arraybackslash}X}
\newcommand{\code}{\nolinkurl}
\newcommand{\src}[1]{\href{run:../#1}{\nolinkurl{#1}}}
\lstset{basicstyle=\small\ttfamily,breaklines=true,columns=fullflexible,
  keepspaces=true,showstringspaces=false,frame=single,
  rulecolor=\color{black!20},backgroundcolor=\color{black!2},
  aboveskip=0.7em,belowskip=0.7em}
\title{UnifiedInterpreter\\HFusion/HIVM Attention 实现说明}
\author{}
\date{2026年10月24日}

\begin{document}
\maketitle
\vspace{-1.5em}

\section{从 attention 的一次执行开始}

当前 CPU 解释器读取归档的 MLIR 文件，为每条操作查找 SSA 操作数，
通过 OCaml effect handler 调用 Rocq 提取的语义函数，最终得到 attention
输出矩阵和每行的 log-sum-exp。两份输入文件为：
\begin{itemize}
  \item \src{artifacts/flash_attention/ir/hfusion_level.mlir}；
  \item \src{artifacts/flash_attention/ir/attention_hivm_compute.mlir}。
\end{itemize}
HFusion 文件包含 \code{linalg}、\code{hfusion} 和公共 dialect。
HIVM 文件使用 \code{hivm.hir.*} 计算原语，并继续使用公共控制流、
张量和内存操作。当前执行原语限定为这两个 attention 截面的子集，
并支持在同一函数中混合使用这些原语。
手工编写的 \src{examples/attention_mixed.mlir} 用于执行这种组合。
仓库中的 H1--H5 演示程序保留在 \src{lib/}，attention 解释器的代码位于
\src{attention/}。

\subsection{输入、输出与计算公式}

固定配置为 batch=head=1，序列长度和 head dimension 均为 128，
采用非 causal attention。输入为
\[
  Q,K,V\in\mathbb{R}^{128\times128},\qquad
  \alpha=1/\sqrt{128}\quad\text{（默认 scale）}.
\]
对每个 query 行 $i$，独立数学参考实现计算
\begin{align}
  S_{ij}&=\alpha\sum_{d=0}^{127} Q_{id}K_{jd},&
  m_i&=\max_{0\le j<128}S_{ij},\\
  E_{ij}&=\exp(S_{ij}-m_i),&
  \ell_i&=\sum_{j=0}^{127}E_{ij},\\
  O_{id}&=\frac{\sum_{j=0}^{127}E_{ij}V_{jd}}{\ell_i},&
  L_i&=m_i+\log\ell_i.
\end{align}
输出 $O$ 为 $128\times128$，输出 $L$ 为 128 项。
这些公式说明参考结果的数学含义；执行器逐条执行输入 MLIR。
参考实现 \code{Attention.reference} 只用于结果比较，
其实现见 \src{attention/attention.ml}。

\subsection{为什么启动 program 0--19}

\code{Attention.run} 顺序启动 20 个逻辑 program，并将编号绑定到
\code{%arg13}。归档 IR 的循环为：
\begin{lstlisting}
scf.for %arg16 = %arg13 to %c4_i32 step %c20_i32 : i32 {
  ...
}
\end{lstlisting}
即从 program id 开始，在小于 4 时执行，步长为 20。因此 program 0--3
分别处理 query 的第 $0$--$31$、$32$--$63$、$64$--$95$、$96$--$127$ 行；
program 4--19 的循环体执行次数为零。编号属于 CPU 解释器的逻辑调用。
同一阶段的 20 次调用共享内存，每次调用使用独立 SSA 环境。
HFusion 与 HIVM 分别从独立内存状态开始执行。

\subsection{一个 query 块经过哪些操作}

以 program 0 为例，\code{memref.reinterpret_cast} 为 $Q$ 的前 32 行建立视图，
\code{memref.copy} 将输入复制到局部 allocation，
\code{bufferization.to_tensor} 建立张量视图。计算过程见表~\ref{tab:attention}。

\begin{table}[htbp]
\centering\small
\begin{tabularx}{\linewidth}{@{}p{0.18\linewidth}YY@{}}
\toprule
计算 & HFusion 文件中的操作 & HIVM 文件中的操作\\
\midrule
$Q_bK^\mathsf{T}$ &
\code{linalg.transpose}、\code{linalg.matmul} &
带 \code{b_transpose} 的 \code{hivm.hir.mmadL1}\\
缩放与逐元素运算 &
\code{linalg.elemwise_binary} &
\code{vmul}、\code{vsub}、\code{vadd}、\code{vdiv}\\
行最大值、行求和 &
\code{linalg.reduce}，执行显式 region &
\code{hivm.hir.vreduce}，由 kind 指定运算\\
扩展行统计量 &
\code{linalg.broadcast} &
reshape 与 \code{hivm.hir.vbrc}\\
指数、对数、转换 &
\code{linalg.elemwise_unary}、\code{hfusion.cast} &
\code{vexp}、\code{vln}、\code{vcast}\\
加权 $V$ 与写回 &
\code{linalg.matmul}、目标写回 &
\code{mmadL1}、目标写回\\
\bottomrule
\end{tabularx}
\caption{归档 attention 中的操作对应关系；HIVM 简写均带 \code{hivm.hir.} 前缀。}
\label{tab:attention}
\end{table}

IR 中还保留了由在线 softmax 形式产生的中间运算，例如由负无穷初值产生的
零权重项。runner 同样执行这些操作。最终通过
\code{bufferization.materialize_in_destination} 写回 $O$ 和 $L$ 的对应行。

\section{代码如何划分职责}

\begin{table}[htbp]
\centering\small
\begin{tabularx}{\linewidth}{@{}p{0.20\linewidth}YY@{}}
\toprule
部分 & 输入与输出 & 主要职责\\
\midrule
xDSL 前端 & MLIR 文本到结构化 JSON & 解析 dialect、类型、属性、SSA、region；记录源码位置\\
OCaml AST 与验证 & JSON 到类型化操作 & 解码字段，检查执行子集、入口签名、形状、属性与作用域\\
OCaml runner & 操作与 SSA 环境到运行值 & 调度函数体、循环、归约 region；构造 effect 请求\\
OCaml handler & 请求到提取函数调用 & 分派计算，适配参数，处理 continuation，提交成功的内存状态\\
Rocq 语义 & 原语参数到结果或内存 computation & 定义张量遍历、广播、归约、矩阵乘法及内存行为\\
数值接口 & 标量操作到 host-float64 结果 & 提供加减乘除、exp/log、NaN 策略和浮点转换约定\\
\bottomrule
\end{tabularx}
\caption{各部分的接口边界。}
\end{table}

Rocq 的 \src{rocq/HFusion/Instr.v} 与 \src{rocq/HIVM/Instr.v}
提供相应语言结构；HIVM 的 \code{request} 定义在
\src{rocq/HIVM/Execution.v}，表示已经求值的原语参数。
文本执行路径使用 \src{attention/parser.ml} 中的操作 AST，
由 runner 转成 handler 请求，再调用提取接口。

\code{rocq/Extraction/Extract.v} 声明导出的语义入口。
Dune 将 Rocq 源文件复制到构建目录，编译后生成
\code{unified_semantics.ml} 和 \code{unified_semantics.mli}，
组成 \code{extracted_semantics} 库。
\code{attention/*.ml}、CLI 和 Python 前端属于手写的执行与适配代码。

\section{原语语义与数值模型}

\subsection{张量值、内存视图与初始化}

\code{DefinedTensor} 携带形状、元素类型和元素列表；
\code{EmptyTensor} 只记录形状与类型，读取元素会返回
\code{ReadUninitialized}。完整覆盖目标的 fill 和广播可以接受空目标；
读取目标初值的矩阵累加需要已初始化张量。

\code{MemRefBackedTensor} 携带 descriptor。
descriptor 包含 allocation 身份、offset、sizes 和 strides。
对于合法逻辑坐标 $\boldsymbol{i}$，寻址形式为
\[
  \operatorname{address}(\boldsymbol{i})
  =\operatorname{offset}+\sum_d i_d\,\operatorname{stride}_d.
\]
\code{to_tensor} 建立视图，后续消费通过 \code{materialize} 读取。
这使多个视图能够引用同一 allocation。定义分别位于
\src{rocq/HIVM/Value.v}、\src{rocq/HFusion/Runtime.v}、
\src{rocq/Memory/Model.v}。

\subsection{HFusion 与 HIVM 的归约}

HFusion 的 \code{sem_linalg_reduce} 使用目标初值和 combiner。
对当前沿第二维归约的矩阵，设目标初值为 $z_i$，则按代码的坐标顺序计算
\[
  r_i=\operatorname{fold}(f,z_i,[x_{i0},\ldots,x_{i,n-1}]).
\]
输出形状为 $[m]$。函数 $f$ 由 MLIR reduction region 的解释执行提供。
定义见 \src{rocq/HFusion/Semantics/Reduce.v}。

HIVM 的 \code{sem_reduce} 输出形状为 $[m,1]$，
从数值接口提供的 identity 开始归约。当前 host-float64 后端为 sum 选择 $0$，
为 max 选择 $-\infty$；目标张量提供形状和类型。
HIVM max 传播 NaN，HFusion region 中的 \code{arith.maxnumf}
在只有一个输入为 NaN 时选择另一个输入。
操作定义与数值实现分别见 \src{rocq/HIVM/Semantics.v}
和 \src{attention/numeric.ml}。

\subsection{矩阵乘法、广播与转换}

HFusion 的 \code{sem_linalg_matmul} 读取目标初值并累加。
HIVM 的 \code{sem_mmad} 根据 \code{mm_init} 选择清零或读取目标，
根据 \code{mm_b_transpose} 选择 $B$ 的行列组织方式。
当前 compute profile 要求 real M/K/N 参数为零，并从张量形状取得实际尺寸。
公共的 \code{matrix_rows}、\code{matrix_columns}、
\code{matrix_product} 在 \src{rocq/Tensor/Matrix.v} 中定义顺序遍历；
每一步乘加调用参数化的标量接口。

当前 HIVM 广播支持标量扩展，以及 $[m,1]\to[m,n]$、广播维度为
$[1]$ 的形式。HFusion attention 使用的广播为 $[m]\to[m,n]$。
操作分别检查输入、目标形状及类型。
HFusion 新增结构操作集中在
\src{rocq/HFusion/Semantics/Structured.v}。

\subsection{host-float64 的具体约定}

当前工作的数值目标是验证解释过程和计算逻辑。
运行时标量使用 OCaml \code{float}；f16/f32 标签继续参与类型检查，
cast 保持浮点数值。exp/log 调用主机数学库；内积按顺序执行乘法和加法。
这些选择定义了当前 CPU 数值模型。

整数操作通过 \src{rocq/HIVM/Arith.v} 处理位宽、符号扩展、
除法和余数；\code{index} 使用该模型选择的 64 位宽度。
OCaml 使用 Zarith 完成与 Rocq 整数表示之间的转换。
程序入口拒绝非有限的 Q/K/V/scale；运算中的 NaN 和无穷值仍按数值接口处理，
结果比较要求全部元素有限。

\section{Effect handler 与 lens}

\subsection{四类 effect 请求}

\src{attention/handlers.ml} 定义四类请求：
\begin{itemize}
  \item \code{Common}：empty、reshape、整数操作及浮点标量原语；
  \item \code{HFusion}：fill、transpose、broadcast、unary、matmul、
    binary、cast、reduce；
  \item \code{HIVM}：调用 \code{Execution.evaluate} 分派的计算 request；
  \item \code{Memory}：调用 \code{run_memory} 执行的内存 computation。
\end{itemize}

以归档 HIVM 文件中的缩放操作为例：
\begin{lstlisting}
%17 = hivm.hir.vmul
  ins(%16, %arg9 : tensor<32x128xf32>, f32)
  outs(%3 : tensor<32x128xf32>) -> tensor<32x128xf32>
\end{lstlisting}
runner 查找 \code{%16}、\code{%arg9} 和 \code{%3}，构造
\code{BinaryRequest (VMul,...)}，通过 \code{Effect.perform}
发出 \code{HIVM} 请求。handler 调用提取的 \code{evaluate}；
\code{sem_binary} 检查形状和类型、扩展标量并逐元素调用
\code{number_binary}。返回张量后，runner 将结果绑定到 \code{%17}。

\code{Effect.Deep.match_with} 安装 handler。
成功时通过 \code{Effect.Deep.continue} 将结果交给后续程序；
语义失败通过 \code{Effect.Deep.discontinue} 传播到请求位置，
由 runner 附加源码上下文。缺少对应 handler 时产生未处理 effect。

\subsection{Rocq computation 与 OCaml effect 的衔接}

Rocq 内存语义使用 \src{rocq/Base/Free.v} 的有限 computation：
\begin{lstlisting}
Inductive computation (E : Type -> Type) (A : Type) : Type :=
| Return (value : A)
| Trigger {X : Type} (request : E X)
    (continue : X -> computation E A).
\end{lstlisting}
\code{Return} 表示结果；\code{Trigger} 保存一次请求及其后续计算。
内存原语构造这样的值，由 Rocq 内存 handler 赋予状态变化。
OCaml 的 \code{Memory} effect 将整段 computation 交给提取的
\code{run_memory} 执行。两处 handler 分别负责 OCaml 请求调度和
Rocq computation 的内存解释。

HFusion reduction 需要额外注意 handler 的动态作用域：
提取的归约函数调用 OCaml combiner，combiner 再执行 MLIR region。
runner 在该回调内部重新安装公共和内存 handler，
并根据 region 内每条操作的阶段选择计算 handler。
\code{arith.maxnumf} 或 \code{arith.addf} 请求使用公共语义。

\subsection{逐操作选择阶段与混合执行}

\src{attention/stage.ml} 定义 \code{Common | HFusion | HIVM}。
OCaml parser 解码操作时填写 \code{operation.stage}：
\code{linalg.*} 和 \code{hfusion.*} 标记为 HFusion，
\code{hivm.*} 标记为 HIVM，支持的公共 dialect 操作标记为 Common。
每个 region 内的操作同样分别标记。
例如 \code{linalg.reduce} 属于 HFusion，
其 region 内的 \code{arith.addf} 属于 Common。

\code{Runner.with_operation_handlers} 根据阶段启用相应的计算 handler，
同时提供公共和内存 handler。外层 \code{with_handlers} 处理公共和内存请求。
验证和执行入口均检查阶段与 opcode 的一致性。
跨阶段执行共用 SSA 环境和内存状态，沿用同一个 memory lens。

混合示例使用 HFusion matmul 产生 \code{%17}，
HIVM vmul 消费它并产生 \code{%18}，随后 HFusion reduce 消费 \code{%18}。
vexp、HFusion cast、vln、输出 vcast 继续交替调用两个阶段的语义。
阶段切换保留原语自身的类型、初始化和归约规则。

\subsection{lens 如何限制状态修改范围}

当前 OCaml 状态为 \texttt{\{memory; marker\}}，
memory lens 读取或替换其中的内存组件。Rocq 定义的状态计算类型为
\[
  M(\mathit{Err},S,A)=S\to\operatorname{result}(\mathit{Err},A\times S).
\]
设 lens 的读写函数分别为 $g:S\to C$ 和 $p:C\to S\to S$，
局部计算为 $c:C\to\operatorname{result}(\mathit{Err},A\times C)$。
\code{zoom} 的定义可写为
\[
\operatorname{zoom}(g,p,c)(s)=
\begin{cases}
 \operatorname{Ok}(a,p(c',s)), & c(g(s))=\operatorname{Ok}(a,c'),\\
 \operatorname{Error}(e), & c(g(s))=\operatorname{Error}(e).
\end{cases}
\]
lawful lens 的三条定律为
\[
g(p(c,s))=c,\qquad p(g(s),s)=s,\qquad
p(c_2,p(c_1,s))=p(c_2,s).
\]
在独立 lens 的假设下，修改内存保持另一组件的观察结果不变。
\code{run_memory_preserves_independent} 证明这一条件性质。
CPU 执行同时检查 \code{marker = 73}，对应具体状态实例的运行检查。

\code{local_memory} 同时检查外层内存错误和内层复合操作错误。
只有成功结果包含可提交的下一状态；OCaml handler 在解包成功结果之后
才更新状态引用。这个失败原子性边界是一次 \code{Memory} 请求；
在此之前已经成功提交的其他请求保持其结果。

\subsection{尾递归内存解释器}

\code{handle_memory_tail} 使用逆序事件累积列表。
\src{rocq/Memory/Handler.v} 中的
\code{handle_memory_tail_correct} 证明：在相同 fuel 下，
它与原 \code{handle_memory} 给出相同结果和内存状态，
事件序列等于初始累积序列恢复顺序后接上原事件序列。
\code{local_memory} 使用空累积列表，因此观察到的事件顺序保持一致。
OCaml 的 primitive trace 另外由 \code{on_call} 记录；
该 JSONL trace 记录提取入口和操作上下文，Rocq 内存事件由内存解释器产生。

\section{解析、验证与控制流}

\code{scripts/attention_frontend.py} 使用固定版本 xDSL 0.74.0
解析整个 module，并调用 \code{module.verify()}。
\code{scripts/attention_dialects.py} 使用 IRDL 和声明式
\code{assembly_format} 定义归档版本所需的操作格式。
公共 dialect 复用 xDSL 的实现。设备元信息作为 opaque attribute 读取。

前端保存原始 SSA 名称以及操作的文件名、偏移和行列，并导出结构化 JSON。
\code{attention/parser.ml} 调用 Python 子进程，用 Yojson 解码为 OCaml AST。
操作记录包含 opcode、stage、输入及目标、操作数类型、结果类型、属性、region、
region 参数与源码位置。源码位置用于报告错误和生成 trace。

\code{Validate.program} 检查入口的 16 个参数、支持的操作和属性、
SSA 引用、局部重复定义、结构操作的形状及类型、reshape 的重关联和 region 参数。
入口绑定中 \code{%arg3}、\code{%arg4}、\code{%arg5} 分别为 Q/K/V，
\code{%arg6} 为 $L$，\code{%arg7} 为 $O$，
\code{%arg9} 为 scale，\code{%arg13} 为 program id。

\code{Runner.execute_block} 持有 SSA 环境并执行操作序列。
\code{scf.for} 使用开区间上界和正步长，
induction variable 的更新调用提取整数加法，并检查更新是否溢出。
每个 region 建立局部参数绑定，因此不同归约 region 可以重复使用局部 SSA 名称。
\code{linalg.yield} 返回 combiner 结果。
执行预算覆盖操作及循环调度；默认 \code{max_steps} 为 $10^6$。

\section{验证证据与适用边界}

\subsection{已经检查的性质}

\begin{longtable}{@{}p{0.23\linewidth}p{0.69\linewidth}@{}}
\toprule
证据 & 检查对象及结论\\
\midrule
\endhead
Rocq 一般性质 &
内存 lens 的独立组件保持，以及尾递归内存 handler 与原定义的等价性。\\
Rocq 泛化示例 &
\code{AttentionStructured.v} 对指定矩阵形状、任意标量和数值接口证明遍历结果相等；
检查 transpose、broadcast、fill。矩阵等式范围为 $2\times3$ 乘 $3\times2$。\\
前端测试 &
真实归档文件、源码位置、局部 SSA、属性与形状拒绝、执行预算及循环 trace。\\
handler 测试 &
真实 float64 矩阵结果、缺少各类 handler、NaN 规则、未初始化和越界读取、
非法 cast，以及复合请求中途失败时不提交状态。\\
端到端测试 &
固定种子、零 Q/K、常量矩阵三组输入，分别解释两个阶段并与独立数学参考比较。\\
修改操作的测试 &
单独的 MLIR 测试文件将 \code{%37} 的 vdiv 改为 vadd。
前端测试检查可执行 AST 的这项差异；执行测试确认 trace 与输出随之改变。\\
混合执行测试 &
\code{mixed_test.ml} 检查逐操作阶段、handler 选择、跨阶段张量传递、
归约回调和共享内存；\code{mixed_execution_test.ml} 比较混合 attention 与独立参考。\\
\bottomrule
\end{longtable}

2026年10月22日的完整验证记录显示：三组输入在两个阶段的 $O$ 和 $L$
最大绝对参考误差均为零，阶段间比较通过。
比较判定使用
\[
  |x-y|\le 10^{-10}+10^{-9}|x|,
\]
并要求两侧结果有限。2026年10月23日编写本文时，重新读取了当前代码，
并运行 \code{scripts/check_attention_output.py} 检查保存的结果与 trace。
该检查通过输出形状、有限性、参考比较结果、program id 覆盖及调用计数检查。

2026年10月24日，完整混合 attention 的测试与 CLI 执行均通过独立参考比较，
$O$ 和 $L$ 的最大绝对误差均为零。逐操作 trace 确认同一次执行调用了
HFusion、HIVM 和公共语义；结果文件检查同时验证形状、program 覆盖和调用计数。
小规模混合程序通过循环、归约回调和内存读写得到预期结果 $[26,26]$。
同日完整执行 \code{dune runtest} 通过，原有三组输入在两个阶段的
$O$ 和 $L$ 参考误差仍均为零。

\begin{table}[htbp]
\centering\small
\begin{tabular}{@{}lrr@{}}
\toprule
默认输入运行统计 & HFusion & HIVM\\
\midrule
runner steps & 66076 & 592\\
primitive calls & 33286 & 570\\
$O$ 最大绝对参考误差 & 0 & 0\\
$L$ 最大绝对参考误差 & 0 & 0\\
\bottomrule
\end{tabular}
\caption{统计来自保存的 CLI 结果；HFusion 显式执行归约 region 中的标量操作。}
\end{table}

默认输入由种子 \code{[2026; 1013]} 生成。
两个阶段得到
\[
 O_{0,0}=0.004908662752173827,\qquad
 L_0=4.853369805731775.
\]
trace 覆盖 program 0--19，只有 0--3 发出 \code{sem_memref_copy} 请求。
测试用例之间通过独立 OCaml domain 并行运行，每个用例内部仍按编号顺序启动 program。

\subsection{结论可以使用到什么范围}

Rocq 原语结构、所列证明和指定输入的数值一致性分别提供相应证据。
当前信任边界包含 xDSL、Python 到 OCaml 的数据转换、手写 runner、
handler 适配、主机数值接口、Rocq 提取及 OCaml 工具链。
完整文本程序的解释正确性、HFusion 到 HIVM 的一般语义保持关系、
低精度数值模型和 NPU 硬件一致性仍需独立验证。

内存模型还有明确的适用限制：当前实现拒绝重叠复制，在创建 reinterpret
视图时检查地址范围，将 empty 的非法消费报告为未初始化读取，
restrict/writable 使用项目定义的内存契约。
这些行为的官方来源核查和适用条件见
\href{semantics-audit-report.md}{语义核查报告}。
当前 attention 使用新 allocation 作为复制目标，视图位于合法范围内，
empty 在消费前由相应操作初始化。

官方规则可从仓库的 \src{vendor/ascendnpu-ir/README.md} 进入固定快照。
其中 \code{HIVMVectorOps.td} 明确给出归约维度保留为 1、
min/max 传播 NaN 等要求。该快照版本为
\nolinkurl{22e3d51df4642fb0f74060af36f5c5c80be80d08}。
具体操作的规范符合性需要逐项比对源码前置条件和结果定义。

\section{构建、运行与代码阅读}

\subsection{可复现命令}

在仓库根目录执行。Python 环境使用 \code{.venv-attention}，
OCaml/Rocq 使用 \code{unified-interpreter} opam switch。
\code{ATTENTION_PYTHON} 可替换为 Python 3.12 及以上解释器的完整路径。

\begin{lstlisting}
ATTENTION_PYTHON=python3.12 sh scripts/bootstrap_attention.sh
export TMPDIR="$PWD/_build/task-tmp"
opam exec --switch=unified-interpreter -- dune build @all
opam exec --switch=unified-interpreter -- dune exec attention -- \
  --stage both \
  --output reports/attention.json \
  --trace reports/attention.jsonl
.venv-attention/bin/python -B scripts/check_attention_output.py
opam exec --switch=unified-interpreter -- dune runtest
\end{lstlisting}

\code{--stage} 接受 \code{hfusion}、\code{hivm}、\code{mixed}、\code{both}。
\code{both} 分别运行两个归档文件；\code{mixed} 要求通过
\code{--ir} 指定包含混合操作的单个文件：
\begin{lstlisting}
opam exec --switch=unified-interpreter -- dune exec attention -- \
  --stage mixed --ir examples/attention_mixed.mlir \
  --output reports/attention-mixed.json \
  --trace reports/attention-mixed.jsonl
\end{lstlisting}
\code{--input} 接受 JSON 对象：Q/K/V 各包含 128 行，每行 128 个数，
scale 为可选有限数。\code{--max-steps} 指定执行预算。
输出包含配置、结果矩阵、误差及调用统计。
trace 中 \code{stage} 表示输入执行模式，\code{ir_stage} 表示操作阶段；
宿主输入准备及结果读取的 \code{ir_stage} 为 null。
其他字段包含 program id、循环变量、源码位置、opcode、结果 SSA 和提取入口。

构建提取库时，\src{extracted/dune} 依赖 Rocq 源码及
\src{rocq/_CoqProject}，调用 \src{scripts/build_extraction.sh}，
在构建目录内完成 Rocq 编译和提取。
因此语义源文件修改可以触发接口重新生成。

\subsection{按验证问题阅读源文件}

以下链接均相对于本文所在的 \code{docs/} 目录。
\begin{longtable}{@{}p{0.28\linewidth}p{0.64\linewidth}@{}}
\toprule
要核查的问题 & 文件与关键入口\\
\midrule
\endhead
实际输入和输出如何绑定 &
\src{attention/attention.ml}：\code{run}、\code{reference}、\code{compare_array}\\
MLIR 如何变成操作对象 &
\src{scripts/attention_frontend.py}：\code{LocatedParser}、\code{read_program}；
\src{scripts/attention_dialects.py}；\src{attention/parser.ml}\\
输入子集如何验证 &
\src{attention/validate.ml}：\code{program}、\code{structured_shape}\\
控制流和 region 如何执行 &
\src{attention/runner.ml}：\code{execute_block}、\code{execute_op}、
\code{with_handlers}\\
混合阶段如何分派 &
\src{attention/stage.ml}：\code{of_opcode}；
\src{attention/runner.ml}：\code{with_operation_handlers}\\
请求如何到达 Rocq &
\src{attention/handlers.ml}：\code{run}、\code{hfusion}、\code{memory_lens}；
\src{rocq/HIVM/Execution.v}：\code{evaluate}、\code{run_memory}\\
两层的计算定义 &
\src{rocq/HFusion/Semantics/Structured.v}；
\src{rocq/HFusion/Semantics/Reduce.v}；
\src{rocq/HIVM/Semantics.v}；
\src{rocq/Tensor/Matrix.v}\\
内存和 lens 的保证 &
\src{rocq/Base/Lens.v}；
\src{rocq/Base/StateM.v}：\code{zoom}；
\src{rocq/Memory/Handler.v}：\code{handle_memory_tail_correct}\\
数值模型如何实现 &
\src{attention/numeric.ml}：\code{scalar}、\code{hivm}\\
如何重新产生证据 &
\src{attention_test/parser_test.ml}；
\src{attention_test/handler_test.ml}；
\src{attention_test/execution_test.ml}；
\src{rocq/Examples/AttentionStructured.v}\\
导出哪些 OCaml 接口 &
\src{rocq/Extraction/Extract.v}；
\src{extracted/dune}\\
\bottomrule
\end{longtable}

\end{document}
